\documentclass[a4paper,12pt]{article}
\usepackage{WSEIRap}
\usepackage{pgf-umlcd}
\usepackage[table]{xcolor}
\usepackage{tikz}
\usetikzlibrary{positioning}
\renewcommand{\titleLab}{Inteligentny Asystent Śniadań}
\renewcommand{\Author}{Mateusz Wnuk}
\renewcommand{\Students}{~}
\renewcommand{\lab}{Zarządznie projektami IT}
\renewcommand{\groupLab}{Gr-6}
\renewcommand{\Date}{1-10-2024}
\renewcommand{\Rodzaj}{Raport}

\begin{document}
    \RapPage


    \section{Inteligentny Asystent Śniadań}\label{sec:inteligentny-asystent-sniadan}

    \subsection{Opis Projektu}\label{subsec:opis-projektu}

    Inteligentny Asystent Śniadań to nowoczesna aplikacja mobilna stworzona z myślą o ułatwieniu codziennych porannych rutyn.
    Główne cele projektu to:
    \begin{enumerate}
        \item \textbf{Automatyzacja Procesu Przygotowywania Śniadań:} Aplikacja pomoże użytkownikom tworzyć spersonalizowane plany śniadaniowe, uwzględniając ich preferencje dietetyczne, dostępne składniki oraz czas, jaki mogą poświęcić na przygotowanie posiłku.
        \item \textbf{Zarządzanie Zasobami Kuchennymi:} Asystent będzie monitorował stan zapasów kuchennych i sugerował zakupy niezbędnych produktów.
        \item \textbf{Porady Żywieniowe i Zdrowotne:} Aplikacja dostarczy użytkownikom cennych informacji na temat wartości odżywczych posiłków, sugerując zdrowe zamienniki oraz optymalne porcje.
        \item \textbf{Integracja z Urządzeniami IoT:} Asystent będzie mógł współpracować z inteligentnymi urządzeniami kuchennymi, takimi jak inteligentne lodówki, kuchenki czy blendery, aby jeszcze bardziej usprawnić proces przygotowywania posiłków.
        \item \textbf{Interaktywne Przepisy:} Użytkownicy będą mieli dostęp do bazy interaktywnych przepisów, które będą dostosowywane na bieżąco w zależności od dostępnych składników oraz preferencji dietetycznych.
        \item \textbf{Automatyczne Przygotowywanie Śniadań:} Asystent będzie mógł sterować urządzeniami kuchennymi, aby automatycznie przygotować śniadanie na zadaną godzinę.
    \end{enumerate}

    \subsection{Wymagania Projektowe}\label{subsec:wymagania-projektowe}

    \subsubsection{Funkcjonalność}
    \begin{itemize}
        \item \textbf{Tworzenie Planów Śniadaniowych:}
        \begin{itemize}
            \item Personalizacja na podstawie preferencji dietetycznych użytkownika.
            \item Uwzględnianie dostępnych składników.
            \item Opcje dla osób z ograniczeniami dietetycznymi (np.\ weganizm, bezglutenowe).
        \end{itemize}
        \item \textbf{Zarządzanie Zapasami:}
        \begin{itemize}
            \item Monitorowanie stanu zapasów kuchennych.
            \item Automatyczne tworzenie list zakupów.
            \item Integracja z platformami zakupowymi online.
        \end{itemize}
        \item \textbf{Porady Żywieniowe:}
        \begin{itemize}
            \item Informacje na temat wartości odżywczych sugerowanych posiłków.
            \item Sugestie zdrowych zamienników.
            \item Personalizowane porady zdrowotne.
        \end{itemize}
        \item \textbf{Integracja z Urządzeniami IoT:}
        \begin{itemize}
            \item Połączenie z inteligentnymi urządzeniami kuchennymi.
            \item Automatyzacja procesów gotowania (np.\ automatyczne ustawianie temperatury w piekarniku).
        \end{itemize}
        \item \textbf{Interaktywne Przepisy:}
        \begin{itemize}
            \item Baza przepisów z możliwością filtrowania według składników, czasu przygotowania, trudności.
            \item Możliwość dodawania własnych przepisów przez użytkowników lub modyfikacji aktualnie dostępnych pod swoje upodobania.
            \item Wskazówki krok po kroku z możliwością odtwarzania wideo.
        \end{itemize}
        \item \textbf{Automatyczne Przygotowywanie Śniadań:}
        \begin{itemize}
            \item Sterowanie urządzeniami kuchennymi (np.
            ekspres do kawy, toster, piekarnik) w celu automatycznego przygotowania śniadania.
            \item Możliwość zaplanowania przygotowywania śniadania na zadaną godzinę.
            \item Monitorowanie procesu przygotowania i dostosowywanie ustawień w czasie rzeczywistym.
        \end{itemize}
    \end{itemize}

    \subsubsection{Technologia}
    \begin{itemize}
        \item \textbf{Platforma:}
        \begin{itemize}
            \item Aplikacja mobilna dostępna na Android i iOS\@.
            \item Responsywna wersja webowa.
        \end{itemize}
        \item \textbf{Integracja IoT:}
        \begin{itemize}
            \item API do komunikacji z inteligentnymi urządzeniami kuchennymi.
            \item Obsługa protokołów takich jak Zigbee, Z-Wave.
        \end{itemize}
        \item \textbf{Baza Danych:}
        \begin{itemize}
            \item Chmurowa baza danych do przechowywania informacji o użytkownikach, przepisach, zapasach.
            \item Bezpieczne przechowywanie danych osobowych.
        \end{itemize}
        \item \textbf{Interfejs Użytkownika:}
        \begin{itemize}
            \item Intuicyjny i łatwy w obsłudze interfejs.
            \item Personalizacja wyglądu aplikacji przez użytkownika.
        \end{itemize}
        \item \textbf{Bezpieczeństwo:}
        \begin{itemize}
            \item Szyfrowanie danych użytkowników.
            \item Regularne aktualizacje zabezpieczeń.
        \end{itemize}
    \end{itemize}

    \subsubsection{Wydajność}
    \begin{itemize}
        \item \textbf{Szybkość Działania:}
        \begin{itemize}
            \item Szybkie ładowanie aplikacji.
            \item Optymalizacja pod kątem wydajności na różnych urządzeniach.
        \end{itemize}
        \item \textbf{Skalowalność:}
        \begin{itemize}
            \item Możliwość obsługi dużej liczby użytkowników.
            \item Skalowalna infrastruktura serwerowa.
        \end{itemize}
    \end{itemize}

    \subsubsection{Utrzymanie i Wsparcie}
    \begin{itemize}
        \item \textbf{Wsparcie Techniczne:}
        \begin{itemize}
            \item 24/7 wsparcie techniczne dla użytkowników.
            \item Regularne aktualizacje aplikacji.
        \end{itemize}
        \item \textbf{Dokumentacja:}
        \begin{itemize}
            \item Szczegółowa dokumentacja dla użytkowników i deweloperów.
            \item Instrukcje i poradniki wideo.
        \end{itemize}
    \end{itemize}

    \subsubsection{Testowanie}
    \begin{itemize}
        \item \textbf{Testy Funkcjonalne:}
        \begin{itemize}
            \item Testowanie wszystkich funkcji aplikacji.
            \item Testy integracyjne z urządzeniami IoT\@.
        \end{itemize}
        \item \textbf{Testy Użyteczności:}
        \begin{itemize}
            \item Badania użyteczności z udziałem rzeczywistych użytkowników.
            \item Regularne zbieranie feedbacku i wprowadzanie poprawek.
        \end{itemize}
        \item \textbf{Testy Bezpieczeństwa:}
        \begin{itemize}
            \item Testy penetracyjne.
            \item Regularne audyty bezpieczeństwa.
        \end{itemize}
    \end{itemize}


    \section{WBS - Work Breakdown Structure}\label{sec:wbs---work-breakdown-structure}
    \begin{center}
        \begin{tikzpicture}[
            every node/.style={font=\small},
            main/.style={rectangle, draw=blue!50!black, fill=blue!20, rounded corners, thick, font=\bfseries\small, minimum width=2cm, minimum height=0.7cm, align=center},
            left/.style={rectangle, draw=green!50!black, fill=green!10, rounded corners, thick, align=center},
            right/.style={rectangle, draw=orange!70!black, fill=orange!10, rounded corners, thick, align=center},
        ]
            \node[main] (projekt) {Projekt};

            \node[left, below left=0.8cm and 1.8cm of projekt] (planowanie) {Faza planowania};
            \node[right, below right=0.8cm and 1.8cm of projekt] (tworzenie) {Faza tworzenia};

            \node[left, below left=0.8cm and -1.3cm of planowanie] (badanie) {Badanie potrzeb rynku};
            \node[left, below right=0.8cm and -1.3cm of planowanie] (projektowanie) {Projektowanie};

            \node[left, below left=0.8cm and -1.1cm of projektowanie] (wizualny) {Wizualny projekt interfejsu};
            \node[left, below right=0.8cm and -1.1cm of projektowanie] (struktura) {Projekt struktury kodu};

            \node[right, below=0.8cm of tworzenie] (implementacja) {Implementacja założeń projektowych};
            \node[right, below right=3.7cm and -5.5cm of implementacja] (integracja) {Integracja z zewnętrznymi urządzeniami};
            \node[right, below left=2.3cm and -2.5cm of implementacja] (utworzenie) {Utworzenie aplikacji};

            \draw[thick] (projekt) -- (planowanie);
            \draw[thick] (projekt) -- (tworzenie);

            \draw[thick] (planowanie) -- (badanie);
            \draw[thick] (planowanie) -- (projektowanie);
            \draw[thick] (projektowanie) -- (wizualny);
            \draw[thick] (projektowanie) -- (struktura);

            \draw[thick] (tworzenie) -- (implementacja);
            \draw[thick] (implementacja) -- (integracja);
            \draw[thick] (implementacja) -- (utworzenie);

        \end{tikzpicture}
    \end{center}

    \begin{tikzpicture}
        \draw (0,0) -- (15,0);

        \foreach \x in {0,3,6,9,12,15}
        \draw (\x,0.15) -- (\x,-0.4);

        \node[anchor=south] at (0,0.15) {Badanie rynku};
        \node[anchor=south] at (3,0.15) {Projekt intefejsu};
        \node[anchor=south] at (6,0.15) {Projekt kodu};
        \node[anchor=south] at (9,0.15) {Pisanie kodu};
        \node[anchor=south] at (12,0.15) {Integracja};
        \node[anchor=south] at (15,0.15) {Gotowy projekt};

    \end{tikzpicture}

    \newpage

    \section{Work Packages - Podejście Waterfall}\label{sec:work-packages---podejscie-waterfall}

    \begin{table}[ht]
        \centering
        \renewcommand{\arraystretch}{1.4}
        \rowcolors{2}{white}{gray!10}
        \begin{tabular}{|>{\columncolor{yellow!20}}c|>{\columncolor{yellow!10}}m{3.2cm}|m{4cm}|m{4.5cm}|>{\columncolor{yellow!10}}m{0.8cm}|}
            \hline
            \rowcolor{yellow!30}
            WP & Nazwa etapu & Cel zadania & Rezultaty (Deliverables) & Czas [dni] \\
            \hline
            WP1 & Badanie rynku & Zidentyfikowanie potrzeb użytkowników & Lista podpunktów do pokrycia przez projekt & 4 \\
            \hline
            WP2 & Projekt interfejsu & Przełożenie pomysłu na warstwę wizualną projektu & Gotowy projekt warstwy wizualnej & 3 \\
            \hline
            WP3 & Projekt kodu & Zaprojektowanie warstwy logicznej aplikacji oraz dobór odpowiednich technologii & Gotowa lista zadań do zrobienia oraz technologie potrzebne do ich zaimplementowania & 3 \\
            \hline
            WP4 & Pisanie kodu & Utworzenie działającej aplikacji & Działająca aplikacja & 7 \\
            \hline
            WP5 & Integracja & Zintegrowanie aplikacji z zewnętrznymi urządzeniami oraz innymi wymaganymi aplikacjami & API używane do integracji zewnętrznych & 3 \\
            \hline
            WP6 & Gotowy projekt & Oddanie gotowego projektu & Koniec prac i pieniądze ze sprzedaży aplikacji & 1 \\
            \hline
        \end{tabular}
    \end{table}


    \section{Wykres Gantta}\label{sec:wykres-gantta}

    \vspace{1cm}

    \begin{ganttchart}[
        hgrid,
        vgrid,
        bar height=0.6,
        bar/. append style={fill=orange!60},
        title label font=\bfseries\footnotesize
    ]{1}{21}
        \gantttitle{ Sprinty Kwiecien 2025}{21} \\
        \gantttitlelist{1,...,21}{1} \\
        \ganttbar{Sprint WP1}{1}{4} \\
        \ganttbar{Sprint WP2}{5}{7} \\
        \ganttbar{Sprint WP3}{8}{10} \\
        \ganttbar{Sprint WP4}{11}{17} \\
        \ganttbar{WP5}{18}{20} \\
        \ganttbar{WP6}{21}{21}
    \end{ganttchart}


    \section{Zasoby projektu}\label{sec:zasoby-projektu}
    Zasoby projektu dzielimy na trzy główne kategorie
    \begin{enumerate}
        \item Zasoby ludzkie
        \item Zasoby technologiczne
        \item Zasoby finansowe
    \end{enumerate}

    \subsection{Zasoby ludzkie}\label{subsec:zasoby-ludzkie}
    W tej kategorii zawierają się wszyscy ludzie współtworzący projekt.
    Programiści, graficy, osoby odpowiedzialne za testy oraz wiele więcej.
    W zależności od projektu może to być zarówno kilka osób jak i naprawdę potężny zespół ludzi.
    Do aktualnego projektu na pewno konieczni będą:
    \begin{enumerate}
        \item Programiści Backend oraz Frontend - odpowiedzialni za działanie aplikacji
        \item UX/UI Designer/s - wygląd oraz wygoda użytkowania gotowego produktu
        \item Electronic engineer - do ewentualnych konsultacji przy integracjach z urządzeniami zewnętrznymi
        \item Testerzy - wyłapywanie błędów i niezgodności z dokumentacją projektu
        \item Project Manager - lider projektu, odpowiedzialny za koordynację całego zespołu
    \end{enumerate}
    Odpowiednie zarządzanie zasobami ludzkimi jest kluczowe dla powodzenia projektu\. Nawet najlepsi specjaliści na nic
    się nie zdadzą jak nie będą wiedzieli co mają robić, albo będą musieli czekać efekty pracy ociągających się kolegów.


    \subsection{Zasoby technologiczne}\label{subsec:zasoby-technologiczne}
    Do tej kategorii można zaliczyć zarówno zasoby używane do pracy nad projektem, jak komputery czy oprogramowanie zainstalowane na nim,
    jak i urządzenia będące składową finalnego produktu, jak telefony i komputery uruchamiające aplikację czy urządzenia zewnętrzne
    przystosowane do zintegrowania z aplikacją mobilną. Do pierwszej kategorii można na pewno zaliczyć:
    \subsubsection{Hardware:}
    \begin{enumerate}
        \item Komputery - na czymś muszą pracować osoby zajmujące się projektem
        \item Rozwiązania chmurowe - umożliwiają łatwą współpracę w zespole, nawet z ludźmi mieszkającymi bardzo daleko od siebie
    \end{enumerate}
    \subsubsection{Software:}
    \begin{enumerate}
        \item Środowiska IDE - aplikacje służące do tworzenia aplikacji(InteliJ IDEA, PHPStorm, VSCode)
        \item Oprogramowanie dla grafików oraz UX/UI Designer'ów(Canva, Adobe Photoshop)
        \item Aplikacje do zdalnej współpracy(Git, Github, Jira, Trello)
    \end{enumerate}
    W drugiej kategorii znajdują się:
    \subsubsection{Hardware:}
    \begin{enumerate}
        \item Wszelkie urządzenia służące do uruchamiania gotowego projektu(dla aplikacji będzie to albo telefon albo komputer/laptop)
    \end{enumerate}
    \subsubsection{Software:}
    \begin{enumerate}
        \item Programy wymagane do uruchomienia aplikacji(dla aplikacji webowej będzie to np. przeglądarka WWW)
    \end{enumerate}

    \subsection{Zasoby finansowe}\label{subsec:zasoby-finansowe}
    Jednak bez pieniędzy żaden projekt nie będzie w stanie ruszyć z miejsca. Komputery do pracy nie należą do tanich,
    tak samo jak oprogramowanie na nich zainstalowane. Również specjaliści w swoich branżach cenią swój czas\. Więc,
    by uzyskać dobrej jakości produkt, trzeba po pierwsze mieć pieniądze, a po drugie mądrze nimi zarządzać.

    \newpage

\section{Koszty}\label{sec:koszty}
    \subsection{Tabela kosztów}\label{subsec:tabela-kosztow}
    \begin{table}[ht]
        \centering
        \renewcommand{\arraystretch}{1.4}
        \rowcolors{2}{white}{gray!10}
        \begin{tabular}{|m{8.2cm}|m{3.4cm}|m{1.4cm}|m{2cm}|}
            \hline
            \rowcolor{yellow!30}
            Zasób & Kategoria & Liczba & Szacowany koszt(zł) \\
            \hline
            Programiści Frontend & Zasoby ludzkie & 2 & 5600 \\
            \hline
            Programiści Backend & Zasoby ludzkie & 3 & 8400 \\
            \hline
            Ux/UI Designer & Zasoby ludzkie & 1 & 1080 \\
            \hline
            Testerzy & Zasoby ludzkie & 3 & 5880 \\
            \hline
            Project Manager & Zasoby ludzkie & 1 & 8400 \\
            \hline
            Komputery dla pracowników & Zasoby techniczne & 10 & 80000 \\
            \hline
            Oprogramowanie dla programistów Frontend & Zasoby techniczne & 2 & 1100 \\
            \hline
            Oprogramowanie dla programistów Backend & Zasoby techniczne & 3 & 1500 \\
            \hline
            Oprogramowanie dla Testerów & Zasoby techniczne & 3 & 600 \\
            \hline
            Oprogramowanie dla Project Managera & Zasoby techniczne & 1 & 300 \\
            \hline
            Aplikacje chmurowe dla zespołu & Zasoby techniczne & 5 & 2500 \\
            \hline
            \multicolumn{3}{|c|}{Podsumowanie kosztów} & 115360 \\
            \hline
            \multicolumn{3}{|c|}{Po doliczeniu marginesu bezpieczeństwa 10\%} & 126896 \\
            \hline

        \end{tabular}
    \end{table}

    \subsection{Rozkład kosztów w czasie}\label{subsec:rozkad-kosztow-w-czasie}
    \begin{ganttchart}[
        hgrid,
        vgrid,
        bar height=0.6,
        bar/. append style={fill=orange!60},
        title label font=\bfseries\footnotesize
    ]{1}{21}
        \gantttitle{ Sprinty Kwiecien 2025}{21} \\
        \gantttitlelist{1,...,21}{1} \\
        \ganttbar{Programiści Frontend}{8}{14} \\
        \ganttbar{Programiści Backend}{8}{14} \\
        \ganttbar{UX/UI Designer}{5}{7} \\
        \ganttbar{Testerzy}{8}{14} \\
        \ganttbar{Project Manager}{1}{21} \\
        \ganttbar{Komputery}{1}{1} \\
        \ganttbar{Oprogramowanie Frontend}{8}{8} \\
        \ganttbar{Oprogramowanie Backend}{8}{8} \\
        \ganttbar{Oprogramowanie Testerzy}{8}{8} \\
        \ganttbar{Oprogramowanie PM}{1}{1} \\
        \ganttbar{Aplikacje chmurowe}{1}{21}
    \end{ganttchart}
\end{document}